\chapter{}
\section{均方差损失下含标签噪声的随机梯度的方差}\label{app:pro1}
\begin{proof}
为了推导和表述的简洁性，我们将
$$
\mathbb{E}_{\zetanoi_n\sim \Dnoi_n}[\Vert \gb_n\left(\omegab,\zetanoi_n\right) - \nabla \fnoi_n(\omegab) \Vert^2 ]
$$
转换为
$$
\mathbb{E}_\epsilon\mathbb{E}_{\zetacle_n}[\Vert \gb_n\left(\omegab,\zetanoi_n\right) - \nabla \fnoi_n(\omegab) \Vert^2 ],
$$
其中 $\epsilon\sim \mathcal{N}(\mu, \sigma^2)$ 表示任意噪声，$\zetacle_n\sim\Dcle_n$。这种转换是等价的，因为任意数据样本中的随机噪声 $\epsilon$ 与干净数据 $(x,y)$ 是相互独立的。同时我们记 $|\zetanoi_n|=|\zetacle_n|=|\zeta_n|$ 且 $|\Dnoi_n|=|\Dcle_n|=|\D_n|$。

根据公式 (\ref{Gradient of clean loss})，我们有
\begin{multline}\label{eq:proof-v1}
\mathbb{E}_\epsilon\mathbb{E}_{\zetacle_n}\left[\left\Vert \gb_n\left(\omegab,\zetanoi_n\right) - \nabla \fnoi_n\left(\omegab\right) \right\Vert^2 \right]
\\=\mathbb{E}_\epsilon\mathbb{E}_{\zetacle_n}\Biggl[\Big\Vert \gb_n\left(\omegab,\zetacle_n\right) -\frac{1}{|\zeta_n|}\sum_{\left(x_i,y_i+\epsilon_i\right)\in \zetanoi_n}2\epsilon_i\nabla h\left(x_i;\omegab\right)
\\ -\left(\nabla \fcle_n\left(\omegab\right)-\frac{1}{|\D_n|}\sum_{\left(x_i,y_i+\epsilon_i\right)\in \Dnoi_n}2\epsilon_i\nabla h\left(x_i;\omegab\right)\right)\Big\Vert^2\Biggl].
\end{multline}

令 $L_n^{noi}\triangleq\frac{1}{|\D_n|}\sum_{\left(x_i,y_i+\epsilon_i\right)\in \Dnoi_n}2\epsilon_i\nabla h\left(x_i;\omegab\right) 
-\frac{1}{|\zeta_n|}\sum_{\left(x_i,y_i+\epsilon_i\right)\in \zetanoi_n}2\epsilon_i\nabla h\left(x_i;\omegab\right)$。则公式 (\ref{eq:proof-v1}) 可表示为
\begin{multline}\label{eq:proof-v2}
    \mathbb{E}_\epsilon\mathbb{E}_{\zetacle_n}\biggl[\Vert \gb_n\left(\omegab,\zetacle_n\right)-\nabla \fcle_n\left(\omegab\right)-L_n^{noi}\Vert^2\biggl]\\=\mathbb{E}_{\zetacle_n}\left[\Vert\gb_n\left(\omegab,\zetacle_n\right)-\nabla \fcle_n\left(\omegab\right)\Vert^2\right]
    +\mathbb{E}_\epsilon\mathbb{E}_{\zetacle_n}\left[\Vert L_n^{noi}\Vert^2\right] \\
    +\mathbb{E}_\epsilon\mathbb{E}_{\zetacle_n}\left[2\langle\gb_n\left(\omegab,\zetacle_n\right)-\nabla \fcle_n\left(\omegab\right),L_n^{noi}\rangle\right].
\end{multline}

在公式 (\ref{eq:proof-v2}) 中，第一项根据假设 \ref{ass:gradient-variance} 小于 $\phi^2_n$；第三项由于干净数据与噪声相互独立而等于零；第二项可表示为
\begin{align}\label{eq:proof-v3}
   &\mathbb{E}_\epsilon\mathbb{E}_{\zetacle_n}\left[\Vert L_n^{noi}\Vert^2\right] \notag \\
    =~&\mathbb{E}_\epsilon\left[\left\Vert \frac{1}{|\D_n|}\sum_{\left(x_i,y_i+\epsilon_i\right)\in \Dnoi_n}2\epsilon_i\nabla h\left(x_i;\omegab\right)\right\Vert^2\right]\notag \\
    &+\mathbb{E}_\epsilon\mathbb{E}_{\zetacle_n}\left[\left\Vert\frac{1}{|\zeta_n|}\sum_{\left(x_{j},y_j+\epsilon_j\right)\in \zetanoi_n}2\epsilon_j\nabla h\left(x_j;\omegab\right)\right\Vert^2\right]\notag \\
    &+\mathbb{E}_\epsilon\mathbb{E}_{\zetacle_n}\biggl[ \frac{1}{|\D_n|}\frac{1}{|\zeta_n|}\sum_{\left(x_i,y_i+\epsilon_i\right)\in \Dnoi_n}\sum_{\left(x_j,y_j+\epsilon_j\right)\in \zetanoi_n}\notag \\
    &4\epsilon_i\epsilon_j\nabla h\left(x_i;\omegab\right)\nabla h\left(x_j;\omegab\right)\biggl].
\end{align}

在公式 (\ref{eq:proof-v3}) 中，基于小批量随机采样，第一项与最后一项之和等于
$$
\frac{4\sigma^2}{|\mathcal{D}_n^{noi}|^2}\sum_{\left(x_{i},y_{i}+\epsilon_{i}\right)\in \mathcal{D}_n^{noi}}\left(\nabla_\omega  h\left(x_{i};\boldsymbol{\omega}\right)\right)^2,
$$
中间项等于
$$
\frac{4\sigma^2}{|\mathcal{D}_n^{noi}||\zeta_n^{noi}|}\sum_{\left(x_{i},y_{i}+\epsilon_{i}\right)\in \mathcal{D}_n^{noi}}\left(\nabla_\omega  h\left(x_{i};\boldsymbol{\omega}\right)\right)^2.
$$
将公式 (\ref{eq:proof-v3}) 代入 (\ref{eq:proof-v2})，即完成证明。若对任意 $\left(x_j,y_j+\epsilon_j\right)\in\D_n$ 都有 $\epsilon_j=0$，该分析依然成立。
\end{proof}
\section{均方差损失下含标签噪声的随机梯度的平方的期望}\label{app:pro2}
\begin{proof}
为简化推导，我们将
$$
\mathbb{E}_{\zetanoi_n\sim \Dnoi_n}[\Vert \gb_n(\omegab, \zetanoi_n) \Vert^2]
$$
表示为 $\mathbb{E}_\epsilon\mathbb{E}_{\zetacle_n}[\Vert \gb_n(\omegab, \zetanoi_n) \Vert^2]$，其中 $\epsilon\sim \mathcal{N}(\mu, \sigma^2)$ 表示任意噪声，$\zetacle_n\sim\Dcle_n$，此为等价变换。

根据公式 (\ref{Gradient of clean loss})，我们有
\begin{multline}\label{eq:proof-s1}
\mathbb{E}_\epsilon\mathbb{E}_{\zetacle_n}\left[\left\Vert \gb_n\left(\omegab, \zetanoi_n\right) \right\Vert^2\right] 
= \mathbb{E}_\epsilon\mathbb{E}_{\zetacle_n}\left[\Vert\gb_n\left(\omegab,\zetacle_n\right)\Vert^2\right] \\
-\mathbb{E}_\epsilon\mathbb{E}_{\zetacle_n}\left[\left\langle\gb_n\left(\omegab,\zetacle_n\right),\frac{4}{|\zeta_n|}\sum_{\left(x_i,y_i+\epsilon_i\right)\in \zetanoi_n}\epsilon_i\nabla h\left(x_{i};\omegab\right)\right\rangle\right]\\
+\mathbb{E}_\epsilon\mathbb{E}_{\zetacle_n}\left[\left\Vert \frac{2}{|\zeta_n|}\sum_{\left(x_i,y_i+\epsilon_i\right)\in \zetanoi_n}\epsilon_i\nabla h\left(x_{i};\omegab\right)\right\Vert^2\right].
\end{multline}

第一项基于假设 \ref{ass:gradient} 被 $G_n^2$ 所限制；第二项由于噪声 $\epsilon_i$ 与数据 $(x_i,y_i)$ 相互独立而等于零；由于小批量随机采样，第三项等于
$$
\sum_{\zeta_n^{noi}\subset\mathcal{D}_n^{noi}} \frac{4\sigma^2}{|\mathcal{D}_n^{noi}||\zeta_n^{noi}|}\sum_{\left(x_{j},y_{j}+\epsilon_{j}\right)\in \zeta_n^{noi}}\left(\nabla_\omega  h\left(x_{j};\boldsymbol{\omega}\right)\right)^2.
$$
由此完成证明。若对任意 $\left(x_j,y_j+\epsilon_j\right)\in\D_n$ 都有 $\epsilon_j=0$，该分析依然成立。
\end{proof}